\chapter{What a block contains}\label{ch:blocks}

This reference describes the protocol that carries a dump. The meaning of the data inside a
block is outside its scope, with one exception worth stating: every block opens with a
fixed signature, and two of them carry readable names. Together these let a program check
a stored dump before sending it back, and list what is in it, without understanding the
rest.

\section{Signatures}\label{sec:signatures}
The bytes below are the first bytes of the reassembled block --- after the payload has been
rejoined from its nibbles, not as they appear on the wire. They were read out of a dump
from a real instrument, and the two rows of each entry are the same sixteen bytes shown
twice, as hexadecimal and as text.

%% GENERATED by notes/sysex-probes/sysex_block_signatures.py --tex
%% Regenerate after any change to the capture or the reassembler

\begin{center}
{\small
\begin{tabular}{llll}
\toprule
block & \textsc{adr} & first sixteen bytes & as text \\
\midrule
\textsc{system, part \& midi} 1 & \bytes{40 00 00} & \bytes{5A 5A 01 00 57 41 30 00} & \texttt{ZZ..WA0.} \\
 & & \bytes{00 00 A0 00 00 00 00 00} & \texttt{........} \\
\textsc{system, part \& midi} 2 & \bytes{40 00 20} & \bytes{78 10 20 20 20 20 20 20} & \texttt{x.~~~~~~} \\
 & & \bytes{20 20 20 20 20 20 20 20} & \texttt{~~~~~~~~} \\
\textsc{sound} & \bytes{20 00 00} & \bytes{57 53 41 20 53 4F 55 4E} & \texttt{WSA~SOUN} \\
 & & \bytes{44 20 52 41 4D 20 53 30} & \texttt{D~RAM~S0} \\
\textsc{combination} 1 & \bytes{50 00 00} & \bytes{5A 5A 5A 5A 00 00 57 53} & \texttt{ZZZZ..WS} \\
 & & \bytes{41 31 20 20 01 00 00 02} & \texttt{A1~~....} \\
\textsc{combination} 2 & \bytes{50 06 00} & \bytes{78 10 20 53 70 61 63 65} & \texttt{x.~Space} \\
 & & \bytes{20 49 6E 66 69 6E 69 74} & \texttt{~Infinit} \\
\bottomrule
\end{tabular}
}
\end{center}


A file whose first bytes do not match the block it claims to be is not a dump of that
category.

\section{Names}
The \textsc{sound} block continues, immediately after its sixteen-byte signature, with
sixteen-byte names in the instrument's own character set, padded with spaces. The second
\textsc{combination} block and the second \textsc{system, part \& midi} block use a
different framing: a two-byte prefix, then a sixteen-byte name, then two further bytes.

A librarian can therefore show the user what a stored file holds --- ``\texttt{SWEEP PAD}'',
``\texttt{Space Infinity}'' --- without decoding anything else.

\begin{note}{Reassembly is self-checking}
A block reassembled correctly comes out at exactly the length its own header declared. A
wrong nibble order, or a payload taken with its trailing flag and checksum still attached,
does not. Checking the length is a cheap way to confirm a receiver is unpacking correctly
before trusting anything it produces.
\end{note}
